%%%%%%%%%%%%%%%%%%%%%%%%%%%%%%%%%%%%%%%%%%%%%%%%%%%%%%%%%%%%%%%%%%%%%%%%%%%%%%%%%%%%%%%%%%%%%%%%%%%
%
% Primordial Machine's Math Matrices Library
% Copyright (C) 2017-2019 Michael Heilmann
%
% This software is provided 'as-is', without any express or implied warranty.
% In no event will the authors be held liable for any damages arising from the
% use of this software.
%
% Permission is granted to anyone to use this software for any purpose,
% including commercial applications, and to alter it and redistribute it
% freely, subject to the following restrictions:
%
% 1. The origin of this software must not be misrepresented;
%    you must not claim that you wrote the original software.
%    If you use this software in a product, an acknowledgment
%    in the product documentation would be appreciated but is not required.
%
% 2. Altered source versions must be plainly marked as such,
%    and must not be misrepresented as being the original software.
%
% 3. This notice may not be removed or altered from any source distribution.
%
%%%%%%%%%%%%%%%%%%%%%%%%%%%%%%%%%%%%%%%%%%%%%%%%%%%%%%%%%%%%%%%%%%%%%%%%%%%%%%%%%%%%%%%%%%%%%%%%%%%

\documentclass[oneside]{book}

%%%%%%%%%%%%%%%%%%%%%%%%%%%%%%%%%%%%%%%%%%%%%%%%%%%%%%%%%%%%%%%%%%%%%%%%%%%%%%%%%%%%%%%%%%%%%%%%%%%
%
% Primordial Machine's Math Matrices Library
% Copyright (C) 2017-2019 Michael Heilmann
%
% This software is provided 'as-is', without any express or implied warranty.
% In no event will the authors be held liable for any damages arising from the
% use of this software.
%
% Permission is granted to anyone to use this software for any purpose,
% including commercial applications, and to alter it and redistribute it
% freely, subject to the following restrictions:
%
% 1. The origin of this software must not be misrepresented;
%    you must not claim that you wrote the original software.
%    If you use this software in a product, an acknowledgment
%    in the product documentation would be appreciated but is not required.
%
% 2. Altered source versions must be plainly marked as such,
%    and must not be misrepresented as being the original software.
%
% 3. This notice may not be removed or altered from any source distribution.
%
%%%%%%%%%%%%%%%%%%%%%%%%%%%%%%%%%%%%%%%%%%%%%%%%%%%%%%%%%%%%%%%%%%%%%%%%%%%%%%%%%%%%%%%%%%%%%%%%%%%

\usepackage[utf8]{inputenc}
\usepackage{textcomp} % for \textlangle and \textrangle macros
\usepackage{xcolor}
%
\usepackage{datetime}
%
\usepackage{tabularx} % for tabularx environment
\usepackage{multirow}
%
\usepackage{lmodern}
\usepackage{microtype}
%
\usepackage{graphicx}
\usepackage[colorlinks]{hyperref}
\usepackage[margin=0.75in]{geometry}
\usepackage{amsmath,amssymb,amsfonts}

\setlength{\parindent}{0pt} % Disable paragraph indention.

\makeatletter

% "(Get|Set)Author".
% The name of the author.
\def\SetAuthor#1{\gdef\@author{#1}}
\def\@author{\@latex@warning@no@line{No author given}}
\def\GetAuthor{\@author}

% "(Get|Set)Email".
% The email (address) of the author.
\def\SetEmail#1{\gdef\@email{#1}}
\def\@email{\@latex@warning@no@line{No email given}}
\def\GetEmail{\@email}

% "(Get|Set)Organization".
% The organization the library is published by.
\def\SetOrganization#1{\gdef\@organization{#1}}
\def\@organization{\@latex@warning@no@line{No organization given}}
\def\GetOrganization{\@organization}

% "(Get|Set)LibraryName".
% The name of the library.
\def\SetLibraryName#1{\gdef\@libraryName{#1}}
\def\@libraryName{\@latex@warning@no@line{No library name given}}
\def\GetLibraryName{\@libraryName}

% "(Get|Set)LibraryIncludesPath".
% The includes directory path.
% Example: "primordialmachine/arithmetic-functors/$(PlatformTarget.toLower())/$(Configuration.toLower())/includes"
\def\SetLibraryIncludesDirectoryPath#1{\gdef\@libraryIncludesDirectoryPath{#1}}
\def\@libraryIncludesDirectoryPath{\@latex@warning@no@line{No library includes directory path given}}
\def\GetLibraryIncludesDirectoryPath{\@libraryIncludesDirectoryPath}

% "(Get|Set)LibraryIncludeFile".
% The include file name of the library.
% Example: "include.hpp".
\def\SetLibraryIncludeFileName#1{\gdef\@libraryIncludeFileName{#1}}
\def\@libraryIncludeFileName{\@latex@warning@no@line{No library include file name given}}
\def\GetLibraryIncludeFileName{\@libraryIncludeFileName}

% "(Get|Set)LibraryIncludeDirectiveFilePath".
% The library include directive file path.
% Example: "primordialmachine/arithmetic_functors/include.hpp".
\def\SetLibraryIncludeDirectiveFilePath#1{\gdef\@libraryIncludeDirectiveFilePath{#1}}
\def\@libraryIncludeDirectiveFilePath{\@latex@warning@no@line{No library include directive file path}}
\def\GetLibraryIncludeDirectiveFilePath{\@libraryIncludeDirectiveFilePath}

% "(Get|Set)LibraryStaticLibrariesDirectoryPath".
% The static libraries directory path.
% Example: "primordialmachine/arithmetic-functors/$(PlatformTarget.toLower())/$(Configuration.toLower())/libraries".
\def\SetLibraryStaticLibrariesDirectoryPath#1{\gdef\@libraryStaticLibrariesDirectoryPath{#1}}
\def\@libraryStaticLibrariesDirectoryPath{\@latex@warning@no@line{No library static libraries directory path given}}
\def\GetLibraryStaticLibrariesDirectoryPath{\@libraryStaticLibrariesDirectoryPath}

% "(Get|Set)LibraryStaticLibraryFileName.
% The static library file name of the library.
% Example: "arithmetic-functors.hpp".
\def\SetLibraryStaticLibraryFileName#1{\gdef\@libraryStaticLibraryFileName{#1}}
\def\@libraryStaticLibraryFileName{\@latex@warning@no@line{No library static library file name given}}
\def\GetLibraryStaticLibraryFileName{\@libraryStaticLibraryFileName}

% "(Get|Set)LibraryVersion".
% The version of the library.
\def\SetLibraryVersion#1{\gdef\@libraryVersion{#1}}
\def\@libraryVersion{\@latex@warning@no@line{No library version given}}
\def\GetLibraryVersion{\@libraryVersion}

% "(Get|Set)LibraryRepository".
% The repository (url) of the library.
\def\SetLibraryRepository#1{\gdef\@libraryRepository{#1}}
\def\@libraryRepository{\@latex@warning@no@line{No library repository given}}
\def\GetLibraryRepository{\@libraryRepository}

% "(Get|Set)DocumentType".
% The document type of this document.
\def\SetDocumentType#1{\gdef\@documentType{#1}}
\def\@documentType{\@latex@warning@no@line{No document type given}}
\def\GetDocumentType{\@documentType}

\let\tableofcontentsORIG\tableofcontents
\renewcommand\tableofcontents{\tableofcontentsORIG}

% Define "maketitle".
\def\maketitle{%
  \thispagestyle{empty}
  \begin{centering}

  \vspace*{6.0\baselineskip} % White space at the top of the page

  %%%%%%%%%%%%%%%%%%%%%%%%%%%%%%%%
  % TITLE
  %%%%%%%%%%%%%%%%%%%%%%%%%%%%%%%%

  \rule{\textwidth}{1.6pt}\vspace*{-\baselineskip}\vspace*{2pt} % Thick horizontal rule
  \rule{\textwidth}{0.4pt} % Thin horizontal rule

  \vspace{0.75\baselineskip} % Whitespace above the title

  {\normalsize\uppercase{\GetOrganization{}}}\\
  \vspace{1\baselineskip} % For some reason the lines are overlapping.
  {\LARGE\uppercase{\GetLibraryName} \uppercase{\GetLibraryVersion}}

  \vspace{0.75\baselineskip} % Whitespace below the title

  \rule{\textwidth}{0.4pt}\vspace*{-\baselineskip}\vspace{3.2pt} % Thin horizontal rule
  \rule{\textwidth}{1.6pt} % Thick horizontal rule

  \vspace{2\baselineskip} % Whitespace after the title block

  %%%%%%%%%%%%%%%%%%%%%%%%%%%%%%%%%
  % SUBTITLE
  %%%%%%%%%%%%%%%%%%%%%%%%%%%%%%%%%

  \vspace{2\baselineskip} % Whitespace after the title block

  \GetDocumentType % Subtitle or further description

  \vspace*{3\baselineskip} % Whitespace under the subtitle

  %%%%%%%%%%%%%%%%%%%%%%%%%%%%%%%%%
  % EDITORS
  %%%%%%%%%%%%%%%%%%%%%%%%%%%%%%%%%

  \vspace*{1\baselineskip} % Whitespace before editors

  Edited By

  \vspace{0.5\baselineskip} % Whitespace before the editors

  {\scshape\Large \GetAuthor \\
  {\normalsize(\href{mailto:\GetEmail}{\GetEmail})} \\} % Editor list

  \end{centering}

  %%%%%%%%%%%%%%%%%%%%%%%%%%%%%%%%%%
  % Footer
  %%%%%%%%%%%%%%%%%%%%%%%%%%%%%%%%%%
  \vspace*{\fill} % Whitespace before the editors
  \begin{tabularx}{\columnwidth}{m{0.75cm}m{4.75cm}X}
  \multirow{2}{=}{\includegraphics[scale=0.05]{primordialmachine-135x135.png}} & Primordial Machine                                                  & \hfill Created with \LaTeX\\
                                                                               & \href{https://primordialmachine.com}{https://primordialmachine.com} & \hfill Compiled on \today\ at \currenttime
  \end{tabularx}
}

\makeatother


\SetOrganization{Primordial Machine's}
\SetLibraryName{Math Matrices Library}

\SetLibraryIncludeFileName{include.hpp}
\SetLibraryIncludesDirectoryPath{primordialmachine/math-matrices/\newline\$(PlatformTarget.toLower())/\$(Configuration.toLower())/includes}

\SetLibraryIncludeDirectiveFilePath{primordialmachine/math/matrices/include.hpp}

\SetLibraryStaticLibrariesDirectoryPath{primordialmachine/math-matrices/\newline\$(PlatformTarget.toLower())/\$(Configuration.toLower())/libraries}
\SetLibraryStaticLibraryFileName{math-matrices.lib}

\SetLibraryVersion{v1.0}
\SetLibraryRepository{https://github.com/primordialmachine/math-matrices}
\SetAuthor{Michael Heilmann}
\SetEmail{michaelheilmann@primordialmachine.com}

\SetDocumentType{User Manual}

\begin{document}

\frontmatter

\begin{titlepage}
\maketitle
\end{titlepage}

\tableofcontents
\addtocontents{toc}{\protect\thispagestyle{empty}}
\pagenumbering{gobble}

\mainmatter

\chapter{Synopsis}
C++ 17 library providing floating point vectors over $n$ dimensions.
The library is made available publicly on
\href{\GetLibraryRepository}{Github}
under the
\href{\GetLibraryRepository/blob/master/LICENSE}{MIT License}.

\chapter{Limitations and Restrictions}
The library officially only supports Visual Studio 2017 and Windows 10.

\chapter{Introductory example}
\textit{\color{orange}This library does not provide any examples yet.}
%Examples are located in the \href{\GetLibraryRepository/blob/master/examples}{examples} directory.

%%%%%%%%%%%%%%%%%%%%%%%%%%%%%%%%%%%%%%%%%%%%%%%%%%%%%%%%%%%%%%%%%%%%%%%%%%%%%%%%%%%%%%%%%%%%%%%%%%%
%
% Primordial Machine's Matrices Library
% Copyright (C) 2017-2019 Michael Heilmann
%
% This software is provided 'as-is', without any express or implied warranty.
% In no event will the authors be held liable for any damages arising from the
% use of this software.
%
% Permission is granted to anyone to use this software for any purpose,
% including commercial applications, and to alter it and redistribute it
% freely, subject to the following restrictions:
%
% 1. The origin of this software must not be misrepresented;
%    you must not claim that you wrote the original software.
%    If you use this software in a product, an acknowledgment
%    in the product documentation would be appreciated but is not required.
%
% 2. Altered source versions must be plainly marked as such,
%    and must not be misrepresented as being the original software.
%
% 3. This notice may not be removed or altered from any source distribution.
%
%%%%%%%%%%%%%%%%%%%%%%%%%%%%%%%%%%%%%%%%%%%%%%%%%%%%%%%%%%%%%%%%%%%%%%%%%%%%%%%%%%%%%%%%%%%%%%%%%%%

\chapter{Building under Visual Studio 2017}
\begin{enumerate}
\item Open the solution \texttt{solution.sln} in Microsoft Visual Studio 2017.
\item Batch build everything.
\item The folder \texttt{packages} contains the distribution of the library i.e. include files and the
      static libraries for
  \begin{enumerate}
    \item the platform targets \texttt{x86} and \texttt{x64} and
    \item configurations \texttt{Release} and \texttt{Debug}.
  \end{enumerate}
\item Copy the contents of the \verb+packages+ folder into a directory. Let
      \verb+[library home]+ be a placeholder denoting the path by which that folder
      can be referenced from your project.
\item Add
  \begin{enumerate}
    \item the include path
\texttt{[library home]/\GetLibraryIncludesDirectoryPath}
	and
    \item the library path
\texttt{[library home]/\GetLibraryStaticLibrariesDirectoryPath}
    to your project.
\end{enumerate}
\item Link your project with the library \texttt{\GetLibraryStaticLibraryFileName}.
\item Add the include directive \texttt{\#include "{}\GetLibraryIncludeDirectiveFilePath"{}} where appropriate.
\item You can now use the functionality provided by the library.
\end{enumerate}


\chapter{Library Interface Documentation}

\section{\texttt{namespace primordialmachine}}
The namespace this library is adding its declarations/definitions to.
The added namespace elements are documented below.

%%%%%%%%%%%%%%%%%%%%%%%%%%%%%%%%%%%%%%%%%%%%%%%%%%%%%%%%%%%%%%%%%%%%%%%%%%%%%%%%%%%%%%%%%%%%%%%%%%%
%
% Primordial Machine's Matrices Library
% Copyright (C) 2017-2019 Michael Heilmann
%
% This software is provided 'as-is', without any express or implied warranty.
% In no event will the authors be held liable for any damages arising from the
% use of this software.
%
% Permission is granted to anyone to use this software for any purpose,
% including commercial applications, and to alter it and redistribute it
% freely, subject to the following restrictions:
%
% 1. The origin of this software must not be misrepresented;
%    you must not claim that you wrote the original software.
%    If you use this software in a product, an acknowledgment
%    in the product documentation would be appreciated but is not required.
%
% 2. Altered source versions must be plainly marked as such,
%    and must not be misrepresented as being the original software.
%
% 3. This notice may not be removed or altered from any source distribution.
%
%%%%%%%%%%%%%%%%%%%%%%%%%%%%%%%%%%%%%%%%%%%%%%%%%%%%%%%%%%%%%%%%%%%%%%%%%%%%%%%%%%%%%%%%%%%%%%%%%%%

%%%%%%%%%%%%%%%%%%%%%%%%%%%%%%%%%%%%%%%%%%%%%%%%%%%%%%%%%%%%%%%%%%%%%%%%%%%%%%%%%%%%%%%%%%%%%%%%%%%
\section{\textit{concept Functor}}
A functor is a struct type providing a constant  \texttt{operator()}.
That operator shall be qualified as \texttt{noexcept}     if possible
and shall be qualified as \texttt{constexpr} if possible.  The return
type of \texttt{operator()} shall be of type \texttt{result}    (or a
cv-qualified variant of that).\newline\noindent{}The functor    shall
provide the member type definition \texttt{result\_type} denoting the
type \texttt{result}.

%%%%%%%%%%%%%%%%%%%%%%%%%%%%%%%%%%%%%%%%%%%%%%%%%%%%%%%%%%%%%%%%%%%%%%%%%%%%%%%%%%%%%%%%%%%%%%%%%%%
\section{\textit{concept BinaryFunctor}}
A \textit{BinaryFunctor} is a \textit{Functor}.
Its \texttt{operator()} has two parameters \texttt{left\_operand} of type
\texttt{left\_operand\_type} (or a cv-qualified variant of that)      and 
\texttt{right\_operand} of type \texttt{right\_operand\_type} (or       a
cv-qualified variant of that).\newline

\noindent{}The functor shall provide the member type              definitions
\texttt{left\_operand\_type}  denoting the type \texttt{left\_operand}    and
\texttt{right\_operand\_type} denoting the type      \texttt{right\_operand}.
\newline

\section{\textit{concept BinaryFunctorBase}}
A  \textit{BinaryFunctorBase} is a template struct type of   name \textit{name}
with three template parameters \texttt{LEFT\_OPERAND}, \texttt{RIGHT\_OPERAND},
and \texttt{ENABLED}. The default value of \texttt{ENALBED} is   \texttt{void}.
Specializations of this template may use \texttt{ENABLED} to    perform SFINAE.
A possible implementation is
\texttt{\newline
\noindent{}template\textlangle typename LEFT\_OPERAND, typename RIGHT\_OPERAND,
typename ENABLED = void\textrangle\newline\noindent{}struct \textit{name};    }
\newline

\noindent{}Its partial specializations are \textit{BinaryFunctors}.
%%%%%%%%%%%%%%%%%%%%%%%%%%%%%%%%%%%%%%%%%%%%%%%%%%%%%%%%%%%%%%%%%%%%%%%%%%%%%%%%%%%%%%%%%%%%%%%%%%%
\section{\textit{concept UnaryFunctor}}
An \textit{UnaryFunctor} is \textit{Functor}.
Its \texttt{operator()} has one parameter \texttt{operand}  of type
\texttt{operand\_type} (or a cv-qualified variant of that).\newline

\noindent{}The functor shall provide the member type definition
\texttt{operand\_type} denoting the type      \texttt{operand}.

\section{\textit{concept UnaryFunctorBase}}
An \textit{UnaryFunctorBase} is a template struct type of     name \textit{name}
with two  template parameters \texttt{OPERAND} and \texttt{ENABLED}. The default
value of \texttt{ENALBED} is \texttt{void}. Specializations of this template may
use \texttt{ENABLED} to perform SFINAE.
A possible implementation is
\texttt{\newline
\noindent{}template\textlangle typename OPERAND,
typename ENABLED = void\textrangle\newline\noindent{}struct \textit{name};    }
\newline

\noindent{}Its partial specializations are \textit{UnaryFunctors}.

%%%%%%%%%%%%%%%%%%%%%%%%%%%%%%%%%%%%%%%%%%%%%%%%%%%%%%%%%%%%%%%%%%%%%%%%%%%%%%%%%%%%%%%%%%%%%%%%%%%
%
% Primordial Machine's Math Matrices Library
% Copyright (C) 2017-2019 Michael Heilmann
%
% This software is provided 'as-is', without any express or implied warranty.
% In no event will the authors be held liable for any damages arising from the
% use of this software.
%
% Permission is granted to anyone to use this software for any purpose,
% including commercial applications, and to alter it and redistribute it
% freely, subject to the following restrictions:
%
% 1. The origin of this software must not be misrepresented;
%    you must not claim that you wrote the original software.
%    If you use this software in a product, an acknowledgment
%    in the product documentation would be appreciated but is not required.
%
% 2. Altered source versions must be plainly marked as such,
%    and must not be misrepresented as being the original software.
%
% 3. This notice may not be removed or altered from any source distribution.
%
%%%%%%%%%%%%%%%%%%%%%%%%%%%%%%%%%%%%%%%%%%%%%%%%%%%%%%%%%%%%%%%%%%%%%%%%%%%%%%%%%%%%%%%%%%%%%%%%%%%

\section{\textit{concept $n,m$-ary-Functor}}
A \textit{$n,m$-ary-Functor} is a \textit{Functor} which provides an constant \texttt{operator()} with $n+m, m, n \geq 0$ parameters.\newline

\noindent{}The first $n$ parameters are inputs to the computation.
The remaining $m$ parameters are parameters to the computation.
The types of all parameters are determined by template parameters.

%\begin{example}
\noindent{}The following example provides a partial specialization of an \texttt{equal\_to\_functor} which is a $2,1$ functor.
The functor determines wether two \texttt{std::array} values are are equal by testing wether the corresponding elements
are equal. The first $2$ parameters of the functor determine the inputs to the computation (2-dimensional vectors) and
the last parameter determines how the elements are compared.
\begin{verbatim}
// Base of 2,1 functor.
template<typename LEFT_OPERAND, typename RIGHT_OPERAND, typename EQUAL_TO, typename ENABLED = void>
struct equal_to_functor;

// 2,1 functor
template<size_t LEFT_ELEMENT, size_t LEFT_NUMBER_OF_ELEMENTS,
         size_t RIGHT_ELEMENT, size_t RIGHT_NUMBER_OF_ELEMENTS,
         typename EQUAL_TO>
struct equal_to_functor<std::array<LEFT_ELEMENT, LEFT_NUMBER_OF_ELEMENTS>,
                        std::array<RIGHT_ELEMENT, RIGHT_NUMBER_OF_ELEMENTS>,
                        EQUAL_TO,
                        std::enable_if_t<std::is_same_v<LEFT_ELEMENT, RIGHT_ELEMENT> &&
                                         LEFT_NUMBER_OF_ELEMENTS == 2                &&
                                         RIGHT_NUMBER_OF_ELEMENTS == 2>>
{
   using left_operand_type = std::array<LEFT_ELEMENT, LEFT_DIMENSIONALITY>;
   using right_operand_type = std::array<RIGHT_ELEMENT, RIGHT_DIMENSIONALITY>;
   using equal_to_type = EQUAL_TO;
   auto operator(const left_operand_type& u, const right_operand_type& v, equal_to_type equal_to) const
   {
     return equal_to(u[0], v[0])
         && equal_to(u[1], v[1]);
   }
};
\end{verbatim}
%\end{example}

\noindent{}The provided \texttt{operator()} shall be qualified as \texttt{noexcept} and \texttt{constexpr} if possible.\newline

\noindent{}$n,m$-functors are frequently (partial) specializations of a base functor.
The base functor shall provide a \texttt{ENABLED} template parameter with a default value \texttt{void} which used to perform SFINAE.
\noindent{}The return value of the \texttt{operator()} represents                the operation
$f : V_1 \times \ldots \times V_n \times P_1 \times \ldots \times P_m \rightarrow(partial) R$.
where $V_i$ is the type of the $i$-th input and $P_j$ the type of the $j$-th parameter. $R$ is
the result of the computation.

\subsection{\textit{concept $m$-BinaryFunctor}}
An \textit{$m$-BinaryFunctor} is a \textit{$2,m$-Ary-Functor}.

\subsection{\textit{concept $m$-UnaryFunctor}}
An \textit{$m$-UnaryFunctor} is a \textit{$1,m$-Ary-Functor}.

\subsection{\textit{concept $m$-NullaryFunctor}}
An \textit{$m$-NullaryFunctor} is a \textit{$0,m$-Ary-Functor}.


\section{\texttt{struct matrix\_traits}}
\begin{verbatim}
template<typename ELEMENT_TYPE, size_t NUMBER_OF_COLUMNS, size_t NUMBER_OF_ROWS>
struct matrix_traits;
\end{verbatim}

\section{\texttt{struct matrix}}
\begin{verbatim}
template<typename TRAITS, typename ENABLED = void>
struct matrix;
\end{verbatim}

\section{Members}

\section{Non members}
%%%%%%%%%%%%%%%%%%%%%%%%%%%%%%%%%%%%%%%%%%%%%%%%%%%%%%%%%%%%%%%%%%%%%%%%%%%%%%%%%%%%%%%%%%%%%%%%%%%
\subsubsection{operator+ (binary, matrix matrix)}
Compute the sum of two matrices.
\begin{verbatim}
template<typename A, typename B>
typename binary_plus_functor<matrix<A>, matrix<B>>::result_type
operator+(const matrix<A>& a, const matrix<B>& b)
{ return binary_plus_functor<matrix<A>, matrix<B>>()(a, b); }
\end{verbatim}
This library provides a specializations for \texttt{matrix\textlangle A\textrangle}
and \texttt{matrix\textlangle B\textrangle} where
\begin{enumerate}
	\item \texttt{A::number\_of\_rows} is equal to \texttt{B::number\_of\_rows},
    \item \texttt{B::number\_of\_columns} is equal to \texttt{B::number\_of\_columns}, and
	\item \texttt{A::element\_type} and \texttt{B::element\_type} are floating point types.
\end{enumerate}

\subsubsection{operator+ (unary, matrix)}
Compute the affirmation of a matrix.
\begin{verbatim}
template<typename A>
typename unary_plus_functor<matrix<A>>::result_type
operator+(const matrix<A>& a)
{ return unary_plus_functor<matrix<A>>()(a); }
\end{verbatim}
This library provides a specializations for \texttt{matrix\textlangle A
\textrangle} where  \texttt{A::element\_type} is a floating point type.

\subsubsection{operator+= (matrix matrix)}
Compute the sum of two matrices, assign the result to the first operand.
\begin{verbatim}
template<typename A, typename B>
auto
operator-=(matrix<A>& a, const matrix<B>& b)
{
   a = a + b;
   return a;
}
\end{verbatim}
%%%%%%%%%%%%%%%%%%%%%%%%%%%%%%%%%%%%%%%%%%%%%%%%%%%%%%%%%%%%%%%%%%%%%%%%%%%%%%%%%%%%%%%%%%%%%%%%%%%
\subsubsection{operator- (binary, matrix matrix)}
Compute the difference of two matrices.
\begin{verbatim}
template<typename A, typename B>
typename binary_minus_functor<matrix<A>, matrix<B>>::result_type
operator-(const matrix<A>& a, const matrix<B>& b)
{ return binary_minus_functor<vector<A>, vector<B>>()(a, b); }
\end{verbatim}
This library provides a specializations for \texttt{matrix\textlangle A\textrangle}
and \texttt{vector\textlangle B\textrangle} where
\begin{enumerate}
	\item \texttt{A::number\_of\_rows} is equal to \texttt{B::number\_of\_rows},
	\item \texttt{A::number\_of\_columns} is equal to \texttt{B::number\_of\_columns}, and
	\item \texttt{A::element\_type} and \texttt{B::element\_type} are floating point types.
\end{enumerate}

\subsubsection{operator- (unary, matrix)}
Compute the negation of a matrix.
\begin{verbatim}
template<typename A>
typename unary_minus_functor<matrix<A>>::result_type
operator-(const matrix<A>& a)
{ return unary_minus_functor<matrix<A>>()(a); }
\end{verbatim}
This library provides a specializations for \texttt{matrix\textlangle A
\textrangle} where  \texttt{A::element\_type} is a floating point type.

\subsubsection{operator-= (matrix matrix)}
Compute the difference of two matrices, assign the result to the left operand.
\begin{verbatim}
template<typename A, typename B>
auto&
operator-=(matrix<A>& a, const matrix<B>& b)
{
   a = a - b;
   return a;
}
\end{verbatim}

\subsection{\texttt{struct transpose\_functor}}
An \textit{UnaryFunctorBase} computing the transpose of a matrix.
A possible implementation is given by
\begin{verbatim}
template<typename OPERAND, typename ENABLED = void>
struct transpose_functor;
\end{verbatim}

\noindent{}This library provides \textit{UnaryFunctor} specializations of this \textit{UnaryFunctorBase}
where \texttt{OPERAND} is specializations of \texttt{matrix}. The result type       \texttt{
result\_type} is \texttt{matrix\textlangle matrix\_traits\textlangle OPERAND\_TRAITS::element\_type,\newline
OPERAND\_TRAITS::number\_of\_rows,OPERAND\_TRAITS::number\_of\_columns\textrangle\textrangle}.

\subsection{\texttt{struct trace\_functor}}

\end{document}
